\subsection{Scoped reuse audit: applicability, denominators and source forms}\label{sec:r8-sok-audit}

\subsection{What the nine items mean}
R1--R9 are checks on whether an observation can be reused for a specified comparison. They are not a quality score for a paper, a safety grade, or a judgment that an original study ought to have performed every possible experiment. Bound means that at least one result record in a family supplies the requested binding; Partial retains both supported and unbound parts; NR-R means the inspected source scope does not report the requested item; NR-V means the source/run/version binding could not be verified; NA means the requested inference is outside the scoped claim. The 20-family matrix is a family-level summary, while result-level records remain distinct where mechanisms or test populations differ.

R4 concerns fitting/calibrating an attacker on protected data when adaptation is part of the reuse question. R5 concerns a matched clean/adapted contrast on the same defended test. Thus R5 does not demand a learned-attacker refit for fixed-rate analytical costs, deterministic namespace controls, or non-learning mechanism benchmarks. Applicability is explicit in the matrix: each item has 11 applicable families for R4 and R5. For R4 the distribution is 4 Bound, 1 Partial, 1 NR-R, 5 NR-V, 9 NA; for R5 it is 2 Bound, 2 Partial, 1 NR-R, 6 NR-V, 9 NA. NR-R and NR-V cells identify missing reporting or binding evidence, not failed experiments.

\begin{table*}[t]\centering\footnotesize
\caption{R4, R5 and R8 family dispositions. Reasons are scoped to the item, not judgments of overall study quality.}\label{tab:r8-family-reasons}
\begin{tabularx}{\textwidth}{@{}lcccX@{}}\toprule
Family & R4 & R5 & R8 & Scope note\\\midrule
UI Traces & B & B & U & Offline attribution does not report objective task utility.\\
WRA & -- & -- & P & Proxy visibility and executable-decoy utility are not fully joined.\\
NetEcho & B & P & U & Topic prediction is not downstream task utility.\\
Early Bird & B & P & U & Semantic-cache timing is not end-to-end task utility.\\
Whisper Leak & ? & ? & U & Post-defense fit/test binding is unverified; replay volume is not task utility.\\
SpliceLeak & -- & -- & U & QCP component cost; no task-utility outcome.\\
Time Will Tell & ? & ? & U & Threshold fit/test binding unverified; analytical padding cost is not task utility.\\
I Know What You Said & -- & -- & -- & Zero-copy component control; family-level utility item inapplicable.\\
Speculative decoding & ? & ? & -- & Learned attack present; protected fit and matched test unverified; payload is cost, not utility.\\
Dangol workflow & B & B & U & Matched result is capture-scoped; objective workflow utility is unreported.\\
Cloud traffic benchmark & ? & ? & P & Classifier under reformulation; semantic preservation is asserted, not task-tested.\\
Remote timing & -- & -- & U & Fixed-rate analysis does not bind task utility.\\
TDXRay & -- & -- & B & Vocabulary coverage is component utility, not deployment-task utility.\\
SafeKV & -- & -- & U & Detection and serving performance use separate conditions.\\
AHB/Turing & P & ? & B & Online task accuracy is bound; publication-run fit/test binding remains unverified.\\
KV-Gov & -- & -- & U & Simulation, hardware gap and HMAC costs are distinct outcomes.\\
CachePrune & -- & -- & B & QA/summarization outcomes bind selective cache sharing.\\
OTRO & -- & -- & -- & Mechanism microbenchmark has no scoped task-utility outcome.\\
SparLeak & ? & ? & P & Local ROUGE proxy is not autonomous workflow utility.\\
Shape and Substance & U & U & U & Inspected scope reports analytic padding, not protected-fit/matched-test or task utility.\\\bottomrule
\end{tabularx}\end{table*}
\noindent\footnotesize B=Bound, P=Partial, U=NR-R, ?=NR-V, --=NA.

For R5, the matrix's applicable denominator is 11, with 2 Bound (2/11), 2 Partial, 1 NR-R and 6 NR-V. Under the narrower interpretation that only explicitly learned/calibrated adaptive-attack experiments call for this matched contrast, Shape and Substance may be recoded NA: 2/10, with 2 Partial and 6 NR-V among the remaining applicable entries. This is a denominator sensitivity, not a correction to the matrix. For R8, the broad item is task \emph{or} component utility: 3 Bound, 3 Partial, 11 NR-R, 3 NA (3/17 applicable). If the estimand is restricted to deployment-task utility, TDXRay's vocabulary-coverage/lookup result is out of scope, yielding 2 Bound, 3 Partial, 11 NR-R and 4 NA (2/16). This alternative does not convert NR-R into an experimental failure. CachePrune's joint R7--R9 binding remains a scoped positive example, not evidence of adaptive robustness.

Across the 20 audited families, the R1--R9 column distributions (Bound/Partial/NR-R/NR-V/NA) are: R1 20/0/0/0/0; R2 17/3/0/0/0; R3 17/3/0/0/0; R4 4/1/1/5/9; R5 2/2/1/6/9; R6 18/2/0/0/0; R7 17/1/2/0/0; R8 3/3/11/0/3; R9 11/3/5/0/1. Each column sums to 20. Machine-readable family mapping and rationale are in \path{provenance/r8-review/family-status-overlay.json}; generated from the main-text matrix by \path{provenance/r8-review/build_overlay.py}.

\subsection{Origin/receiver/polarity deductions}
Origin A/B labels are indexing coordinates for where signal is generated; they are not causal classes, a receiver projection, or a portability guarantee. Three record-level deductions illustrate the use of the map. First, NetEcho's upstream packet-size defense (Origin B) cannot be read as a guarantee about an agent's Origin-A tool-call sequence; after plaintext token consumption, a frontend's re-emission creates a new B size trace. Second, WRA's Origin-A read-only decoys protect user intent (Type I), whereas AHB's online actions evade a platform detector (Type III) and can mutate state; shared origin does not imply shared utility direction. Third, identical Origin-A read/write dispatch projections can preserve an inference feature while writes alter the account/database, so state-preservation claims require execution-path and final-state evidence (E6). These correct specific cross-record transfers; they are not new theorems.

\subsection{Citation-form inventory and checked primary pages}
The compiled main bibliography has 77 cited entries: 49 arXiv-only labels; 16 venue labels; 5 venue labels that also cite a preprint; 1 book; 2 standards; and 4 vendor documents. These counts classify current bibliography form only. They do not certify peer review, and arXiv-only does not mean unpublished. The complete row-level list is \path{provenance/r8-review/all-citations-form-labels.csv}; the 20-family subset is \path{provenance/r8-review/publication-form-labels.csv}.

For the eight cited case records, primary-page lookup on 24 October 2026 found: NetEcho's official arXiv page records v1 (\url{https://arxiv.org/abs/2510.25472}); SplitDecoy v1 (\url{https://arxiv.org/abs/2609.04382}); AHB/Turing v1 (\url{https://arxiv.org/abs/2604.09574}); Dangol v3, revised 17 June 2026 (\url{https://arxiv.org/abs/2606.07150}); CachePrune v1 (\url{https://arxiv.org/abs/2605.23640}); and OTRO v3, revised 7 October 2026 (\url{https://arxiv.org/abs/2606.17358}). No venue statement surfaced on those inspected arXiv pages; this bounded lookup does not prove absence of a venue publication. The official USENIX page confirms WRA in USENIX Security 2026 (\url{https://www.usenix.org/conference/usenixsecurity26/presentation/jeong}). IEEE's official result lists Early Bird in TIFS 20 (2025), pp. 11431--11446, publication 17 October 2025 (\url{https://ieeexplore.ieee.org/document/11206380/}); direct page access was blocked by a JavaScript/robot interstitial. Page observations and limitations are preserved in \path{provenance/r8-review/official-source-status-checks.json}.

